% Gauss, Werke X/1 (1917), native81 lower–84, printed74 lower–77.
\begin{center}\large\textsc{Gauss} an \textsc{Olbers}.\end{center}
\HistoricCitation{[\textsc{Wilhelm Olbers}, Sein Leben und seine Werke II, 1 (1900), S. 376, 377.]}
\begin{flushright}Braunschweig, 1807 Juli 21.\end{flushright}
\noindent .\quad .\quad .\quad .\quad .\quad .\par
\par\indent Bei meiner Rückkunft habe ich hier einige Briefe aus Paris vorgefunden
von \textsc{Bouvard}, \textsc{Lagrange} und \textsc{Sophie Germain} .\ .\ .\ .\quad Meinen Brief hat \textsc{Lagrange}
\par\vfill\noindent\rule{.25\linewidth}{.4pt}\par\smallskip
\begingroup\footnotesize\noindent[*) Werke I, S. 107 ist dieser Druckfehler berichtigt.]\par\endgroup
\clearpage
\SourceLeaf{082}\ArithmeticHead{75}{R}{GAUSS AN OLBERS.}
\noindent wirklich, wie er selbst schreibt[*)], sogleich im \textls[100]{Institut} und \textls[100]{Bureau
des Longitudes} vorgelesen .\ .\ .\ .\ .

\textsc{Lagrange} interessirt sich noch mit vieler Wärme für die Astronomie
und höhere Arithmetik; die beiden Probe-Theoreme (in welchen Primzahlen
2 ein kubischer oder ein biquadratischer Rest ist), die ich auch Ihnen vor
einiger Zeit mittheilte, hält er für »ce qu’il peut y avoir de plus beau et de
plus difficile à demontrer«. Aber die \textsc{Sophie Germain} hat mir die Beweise
derselben geschickt; noch habe ich sie zwar nicht durchgehen können, ich
glaube aber, dass sie gut sind; wenigstens hat sie die Sache von der rechten
Seite angegriffen, nur etwas weitläufiger sind sie als nöthig sein wird. .\ .\ .\ .
\ClosingRule
\begin{center}\large\textsc{Gauss} an \textsc{Olbers}.\end{center}
\HistoricCitation{[\textsc{Wilhelm Olbers}, Sein Leben und seine Werke II, 1 (1900), S. 629.]}
\begin{flushright}Göttingen, 1816 März 21.\end{flushright}
\noindent .\quad .\quad .\quad .\quad .\quad .\par
\par\indent Für Ihre Nachrichten, die Pariser Preise betreffend, bin ich Ihnen sehr
verbunden. Ich gestehe zwar, dass das \textsc{Fermat}sche Theorem als isolirter
Satz für mich wenig Interesse hat, denn es lassen sich eine Menge solcher
Sätze leicht aufstellen, die man weder beweisen, noch widerlegen kann. Allein
ich bin doch dadurch veranlasst, einige alte Ideen zu einer \textls[100]{grossen} Erweiterung
der höheren Arithmetik wieder vorzunehmen. Freilich gehört diese
Theorie zu den Dingen, wo man nicht voraussehen kann, inwiefern es gelingen
wird, dunkel vorschwebende entfernte Ziele zu erreichen. Ein glückliches
Gestirn muss mit obwalten, und meine Lage und so vielfache abziehende
Geschäfte erlauben mir freilich nicht, solchen Meditationen so nachzuhängen,
wie in den glücklichen Jahren 1796—1798, wo ich die Hauptsachen meiner
\textit{Disquisitiones Arithmeticae} bildete. Allein ich bin überzeugt, wenn das \textls[100]{Glück}
mehr thun sollte, als ich erwarten darf, und mir einige Hauptschritte in jener
\par\vfill\noindent\rule{.25\linewidth}{.4pt}\par\smallskip
\begingroup\footnotesize
\noindent[*) Der Brief von \textsc{Lagrange} vom 31. Mai 1807 ist abgedruckt \textit{Oeuvres de Lagrange} t. XIV, 1892, S. 298.]\par
\endgroup
\noindent\hfill10*\par\clearpage
\SourceLeaf{083}\ArithmeticHead{76}{L}{ARITHMETIK.\quad NACHLASS.}
\noindent Theorie glücken, auch der \textsc{Fermat}sche Satz nur als eines der am wenigsten
interessanten Corollarien dabei erscheinen wird.\quad .\quad .\quad .\quad .\quad .\quad .\quad .\quad .\quad .\quad .
\ClosingRule
\begin{center}\large\textsc{Gauss} an \textsc{Bessel}.\end{center}
\HistoricCitation{[Briefwechsel zwischen \textsc{Gauss} und \textsc{Bessel}, Leipzig 1880, S. 246.]}
\begin{flushright}Göttingen 23. December 1816.\end{flushright}
\noindent .\quad .\quad .\quad .\quad .\quad .\par
\par\indent Seit mehreren Monaten sind es gewisse Untersuchungen aus der höheren
Arithmetik, auf die ich wiederum zurückgekommen bin, und die mich schon
seit beinahe 12 Jahren geplagt haben. Sie gehören zu der Gattung derjenigen,
wo man nicht im voraus sagen kann: dies will ich thun, sondern wo, vielleicht
nach 999 misslungenen Versuchen, eine glückliche 1000-te Combination
zum Ziele führt. Jetzt habe ich zwar das Ziel erreicht, doch immer noch
auf einem nicht genug kurzen Wege. Der Gegenstand ist die Theorie der
biquadratischen Reste, deren ich vielleicht schon mehrere Male gegen Sie
erwähnt habe. Auch diess Brüten über einer Sache, ohne dass mittheilbare
Resultate daraus hervorgehen, hat mich in anderen Dingen und besonders in
meiner Correspondenz zurückgesetzt. Ich werde jetzt nur soviel davon aufschreiben,
dass die neuen noch in der Luft schwebenden Ideen wenigstens
meinem Gedächtniss erhalten werden.
\begin{center}\textbf{BEMERKUNG.}\end{center}
\begingroup\small
\par\indent Der hier in der ursprünglichen Schreibung wiedergegebene Brief an \textsc{Sophie Germain} ist zuerst 1879
von dem Fürsten B. \textsc{Boncompagni} als Autographie veröffentlicht worden. Die Urschrift befindet sich jetzt
im \textsc{Gauss}archiv, während zwei frühere und ein späterer Brief von \textsc{Gauss} an dieselbe Empfängerin in der
\textls[100]{Bibliothèque nationale} zu Paris aufbewahrt werden. Die Briefe der \textsc{Sophie Germain} an \textsc{Gauss} sind
veröffentlicht unter dem Titel: \textit{Cinq lettres de Sophie Germain à Charles Frédéric Gauss publiées par
B. Boncompagni d’après les originaux possédés par la Société Royale des sciences à Göttingen. Berlin 1880.}
Über die Bedeutung des hier abgedruckten Briefes vom 30. April 1807 vgl. E. \textsc{Schering}, Göttinger Nachrichten
1879, S. 381ff.

Die Nachricht, auf die sich \textsc{Gauss} in dem Schreiben an \textsc{Olbers} vom 21. März 1816 bezieht, ist
in einem Briefe vom 7. März 1816 enthalten, den \textsc{Olbers} aus Bremen an \textsc{Gauss} gerichtet hat*). Sie betrifft
\par\vfill\noindent\rule{.25\linewidth}{.4pt}\par\smallskip
\noindent*) W. \textsc{Olbers}, Sein Leben und seine Werke II, 1, S. 626.\par
\clearpage
\SourceLeaf{084}\ArithmeticHead{77}{R}{BEMERKUNG.}
\noindent eine von der Pariser Akademie für 1818 gestellte Preisaufgabe, von der \textsc{Olbers} schreibt: [sie] »scheint mir
recht für Sie geschaffen, lieber \textsc{Gauss}.« Die Aufgabe lautet, soweit sie hier in Betracht kommt, in der
\textsc{Olbers}schen Übersetzung wie folgt:

».\ .\ . Nun ist also nur noch das andere Theorem zu beweisen: \textls[100]{Daß es nämlich keine Potenz
außer dem Quadrat gibt, die in zwei andere Potenzen von demselben Grade zerlegt
werden könne.}«
\begin{flushright}\textsc{Schlesinger.}\end{flushright}
\ClosingRule\clearpage\endgroup
